\documentclass[10pt,a4paper]{amsart}
\usepackage[margin=19mm]{geometry}
\usepackage{amsmath,amssymb,mathtools}
\usepackage[scr=rsfs,cal=euler]{mathalfa}
\usepackage{xfrac}
\usepackage[unicode,hidelinks,bookmarks=false]{hyperref}
\newcommand{\ie}{i.e.\ }
\newcommand{\cf}{cf.\ }
\newcommand{\resp}{respectively\ }
\newcommand{\Subsection}{\subsection}
\providecommand{\nobreakdash}{\nobreak\hbox{-}}
\setlength{\emergencystretch}{2em}
\multlinegap0pt
\usepackage{amssymb}
\usepackage{tikz}
\newtheorem{lemma}{Lemma}
\title{The Discrete Harmonic Center of a Quadrilateral}
\author{Marc Alexa}
\date{Source: arXiv:2609.00917v1 (CC BY 4.0)}
\begin{document}
\maketitle

\noindent\emph{Source and licence.} The Discrete Harmonic Center of a Quadrilateral. Version arXiv:2609.00917v1 (CC BY 4.0), distributed by arXiv under the Creative Commons Attribution 4.0 International licence, http://creativecommons.org/licenses/by/4.0/, as recorded in the arXiv licence metadata for this exact version and shown on the abstract page https://arxiv.org/abs/2609.00917v1. Full licence notice: \texttt{licenses/arxiv-2609.00917-CC-BY-4.0.txt}.\par\smallskip

\noindent\textit{Editorial scope.} Counted unit: the article's lemma lemma (source0.tex lines 2499--2501). The article's complete original proof of it is delivered with it (source0.tex lines 2502--2522, one \texttt{proof} environment of 21 lines). No mathematics is written, repaired or re-derived: the statement and the proof are the article's own lines, in the article's own notation, and no retained line is substituted or re-flowed.  No further source-local context is needed: every cross-reference of every retained line resolves inside the two delivered spans.  Nothing was omitted from any retained span, so the package carries no omission record.  No retained line contains a citation, so the wrapper carries no bibliography and no bibliography entry.  No retained line contains a figure environment, an image-inclusion command or drawing code, so no image is delivered and the package declares no asset dependency.  Editorial formatting only, in addition: an AMS \texttt{amsart} preamble that restates the article's own declarations and macros used by the retained lines, namely the environments \texttt{lemma} on separate counters, together with the standard packages those lines need.  The retained environments print in this document as Lemma 1 in order of appearance, whereas the article prints the same items with the numbering of its own full document.  The complete native source of the article as distributed in the arXiv e-print for this exact version is bundled unchanged as \texttt{sources/209069/source0.tex}.
\begin{lemma}
    The points $O$, $P$, $z^*$, $\tilde{z}^*$ are on a common line. 
\end{lemma}

\begin{proof}
The two circles $\gamma_{02}$ and $\gamma_{13}$ intersect in $z^*$ and $\tilde z^*$. The line through these intersections is the \emph{radical axis} of the two circles defined by 
\begin{equation}\label{eq:radicalline}
    |x-o_{02}|^2 - r_{02}^2 = |x-o_{13}|^2 - r_{13}^2.
\end{equation}
Because $\Gamma$ is orthogonal to $\gamma_{02},\gamma_{13}$ we have
\begin{equation}
    R^2+r_{02}^2 = |O-o_{02}|^2,
    \qquad
    R^2+r_{13}^2 = |O-o_{13}|^2    
\end{equation}
and using this when plugging $O$ into Eq.~\ref{eq:radicalline} demonstrates that $O$ is on the radical line.\\
Each of the diagonals in the quadrilateral is part of the radical line of  $\Gamma$ and either $\gamma_{02}$ or $\gamma_{13}$, meaning the intersection of the diagonals $P$ satisfies both
\begin{equation}
    |x-o_{02}|^2 - r_{02}^2 = |x-O|^2 - R^2
    \qquad \text{and}
    \qquad
    |x-o_{13}|^2 - r_{13}^2 = |x-O|^2 - R^2.  
\end{equation}
Since the right hand sides are identical, $P$ satisfies Eq.~\ref{eq:radicalline} and lies on the common line. 
\end{proof}

\end{document}
